% This is samplepaper.tex, a sample chapter demonstrating the
% LLNCS macro package for Springer Computer Science proceedings;
% Version 2.20 of 2017/10/04
%
\documentclass{article}
%
\usepackage{subfig}
\usepackage{graphicx}

\usepackage{cite}

\usepackage{todonotes}

\usepackage[utf8]{inputenc}
\usepackage[acronym]{glossaries}

%---------Tree:Begin-----------------
\usepackage{tikz}
%%%<
\usepackage{verbatim}
%%%>
\usetikzlibrary{trees}
%----------Tree:End------------------

%----------Glossary:Start------------
\makeglossaries

% Sample glossary entry
\newglossaryentry{maths}
{
    name=mathematics,
    description={Mathematics is what mathematicians do}
}

% Sample acronym
\newacronym{mt}{MT}{Math Teachers}

%----------Glossary:End--------------

%----------Title:Start---------------
\title{The \LaTeX\  IPA Template}
\author{Max Muster}
\date{1. Mai 2020}
%\institute{Quality \& Quantity GmbH, Musterstrasse 12, 0000 Stadt}
%----------Title:End-----------------

\begin{document}

\maketitle

\newpage
\tableofcontents

%Start of Report
\part{Umfeld und Ablauf}
\section{Zielgruppe}
\Gls{maths} is very Important for \acrfull{mt}
\section{Aufgabenstellung (gemäss PkOrg)}
\begin{itemize}
    \item{}
    \item{}
    \item{}
    \item{}
\end{itemize}
\section{Deklaration der Firmenstandards}
\section{Projektaufbauorganigramm}
\section{Organisation}
\section{Ablauf}
\section{Meilensteine}
\section{Zeitplan}
\section{Arbeitsjournal}
\section{Projektjournal}

\part{Projekt-Dokumentation}
\section{Kurzfassung}
\section{Initialisierung}
\begin{figure}[h!]
    \subfloat[Ich beim Arbeiten]{\label{me_working:a}\includegraphics[width=0.3\columnwidth]{figures/hard_working_1.jpg}}
    \hfill
    \subfloat[Ich beim Arbeiten]{\label{me_working:b}\includegraphics[width=0.3\columnwidth]{figures/hard_working_2.jpg}}
    \hfill
    \subfloat[Ich beim Arbeiten]{\label{me_working:c}\includegraphics[width=0.3\columnwidth]{figures/hard_working_3.jpg}}
    \label{me_working}
    \caption{Drei Beispiele wie ich gearbeitet habe.}
\end{figure}
\section{Konzeptphase}
\section{Realisierung}
\begin{figure}[h!]
    \centering
    \includegraphics[width=0.4\columnwidth]{figures/my_product.jpg}
    \caption{Ein Bild meiner Software}
    \label{my_product}
\end{figure}
\section{Testphase}
\section{Einführung}
\section{Schlusswort}


\newpage

\setcounter{tocdepth}{1}
\listoffigures

\newpage
\printglossary[type=\acronymtype]
\newpage
\printglossary


\end{document}
